\section{When later disclosures change the historical query}
\label{sec:apple-pronouncements}

A later filing can reveal more about an earlier period without identifying a
different economic history. For example, a new segment-expense disclosure may
make a previously hidden cost visible, while the transactions that generated
that cost remain the same. To use that information correctly, a model must
retain the original report, the later presentation and the date at which the
additional detail became available. Apple's announcement and its later
segment tables show the difference between a planned disclosure and the
comparative information eventually reported.

Annual and interim MD\&A describe planned changes to tax and segment disclosures;
the interim adds expense disaggregation and a subsequent effective-date
clarification \citep[Recent Accounting Pronouncements, printed p.~26]{apple2024}
\citep[Recent Accounting Pronouncements, printed p.~17]{apple2025q1}.
The following table preserves Apple's stated adoption plans. The quarters are
fiscal, and the descriptions are issuer reports of the updates, rather than an
independent determination of their operative requirements.
\begin{center}
\small
\begin{tabular}{p{0.19\textwidth}p{0.40\textwidth}p{0.25\textwidth}}
\toprule
Update & Reported disclosure change & Apple's stated plan\\\midrule
ASU 2023-07 & Segment expenses and information about the chief operating
decision maker & Fiscal 2025 Q4, retrospective\\
ASU 2023-09 & Income-tax reconciliation and taxes-paid detail & Fiscal 2026 Q4,
prospective\\
ASU 2024-03, clarified by 2025-01 & Expense components and selling-expense
amount and definition & Fiscal 2028 Q4, prospective\\\bottomrule
\end{tabular}
\end{center}

The tax discussion says the rate reconciliation will include additional
information and detail for items meeting a quantitative threshold. Taxes paid
will be disaggregated into federal, state and foreign taxes, with further detail
for significant individual jurisdictions. It gives neither the threshold value
nor the future jurisdictional amounts. Those cash-tax populations also remain
distinct from tax expense, payable and deferred balances; their reconciliation
is developed in Section~\ref{sec:apple-tax-settlement}.
An announced disclosure change cannot be used to allocate an old tax total
among jurisdictions by assumption.

The earlier segment discussion concerns expenses that are significant and regularly
provided to the chief operating decision maker, together with that person's
title, position and use of segment profit or loss in assessing performance and
allocating resources. That anticipated disclosure could answer questions about
the information used by management. It did not yet identify the regional
allocation of corporate expenses or supply segment assets. The adoption
sentence was a dated plan at that information cutoff. The retained fiscal
2025 annual report now lets us check what followed.

Apple reports that it adopted ASU 2023-07 for its 2025 annual period using a
retrospective transition method \citep[Note 1]{apple2025annual}. Its Note 13
presents 2024 and 2023 comparative segment expenses. Take the Americas in
fiscal 2024. The earlier Note 13 reports net sales of $167{,}045$ million
and operating income of $67{,}656$ million, but no separate regional cost of
sales or selling-and-marketing amount \citep[Note 13, printed p.~47]{apple2024}.
The later comparative gives $89{,}587$ million of cost of sales and $9{,}802$
million of selling and marketing for that same region and fiscal year
\citep[Note 13, printed p.~47]{apple2025annual}. The displayed relationship is
\[
167{,}045-89{,}587-9{,}802=67{,}656\quad\text{(USD millions).}
\]
The 2025 table gives corresponding expense rows for the other four regions.
Together, their disclosed selling-and-marketing expenses total $18{,}639$
million. The five operating incomes still sum to $162{,}044$ million; after
the displayed corporate research and development of $31{,}370$ million and
general and administrative amount of $7{,}458$ million, the company reports
the same $123{,}216$ million of operating income. The 2024 report called its
$7{,}458$ million line ``other corporate expenses, net'' and qualified its
contents. The identical displayed amount is not, by itself, proof that the
two labels have identical scope. Both original and later assertions retain
their source identities.

This later presentation answers a prior-period disclosure question without
identifying new fiscal-2024 transactions or creating a correction journal.
Reading the named 2024 source leaves the regional expense detail unresolved;
reading the later source supplies its 2025 comparative assertions. An
investor-information view additionally needs evidence that the exact chosen
source version was publicly available by its requested cutoff. The accepted
corpus and filing-event metadata do not prove that for the retained 2024 PDF
or the later issuer-hosted HTML. Nor do the new regional expense rows reveal
the cross-classified product-by-region sales, entity cash allocation, or
behavioural response needed by the earlier regional stress example.

The interim expense passage reports ASU 2024-03's issuance in November 2024
and ASU 2025-01's clarification in January 2025. It describes purchases of
inventory, employee compensation, depreciation and intangible-asset amortization,
as applicable, within certain expense captions, a qualitative description of
the remaining amounts, and both the amount and the company's definition of
selling expenses. This describes disclosure populations; it does not provide
their values or make selling expenses identical to the entire selling, general
and administrative line. Purchases included in an expense disclosure likewise
cannot be identified with cash paid to suppliers without the relevant timing
and population facts.

The observed segment comparatives show why the two source-content views must
remain separate: the original cannot be overwritten even when a later filing
reports the same amounts with more detail. This comparison checks the affected
periods, dimensions, displayed units and issuer labels; it does not use the
phrase ``retrospective transition'' as an amount. The prospective tax and
expense plans still supply no backward allocation of missing components.
Their exact transition requirements need the operative updates and any
applicable amendments.

The January clarification illustrates the timing problem directly. It occurs
after the interim period ended on 28 December 2024 but is described in the
interim report. An event can therefore belong to the information in a report
without occurring inside its reporting period. A historical forecast cutoff
needs evidence of availability and the exact retained version, not merely the
period-end or the update's issuance month. Likewise, the earlier announced
segment adoption did not become an observed adoption until the later filing
was checked. The tax and expense plans remain future claims until their own
later evidence is inspected.

\subsection{Unchanged critical estimates do not mean constant amounts}

Annual Critical Accounting Estimates says management makes judgments,
assumptions and estimates affecting reported amounts, using historical
experience and other assumptions it considers reasonable to judge carrying
values. Its uncertain-tax discussion says reserves can be adjusted as facts
and circumstances change, and that final outcomes may differ from those
reserves. Its legal-contingency policy is the incurred-loss and estimability
condition examined above \citep[Critical Accounting Estimates, printed p.~26]{apple2024}.
These are descriptions of judgment and policy, not a disclosed numerical
estimation model or an assurance that the estimates will prove accurate.

The interim references both annual and interim Note 1 and reports no material
changes to critical accounting estimates since the annual filing
\citep[Critical Accounting Estimates, printed p.~17]{apple2025q1}.
That statement preserves the annual dependency and its materiality qualifier.
It cannot make every carrying amount, reserve or underlying assumption exactly
constant, and it supplies no numerical tolerance for a comparison. A model
should retain the reported policy and qualified continuity assertion alongside
the actual dated balances, then request the missing estimate or transition
facts when the query needs them. These estimate passages supply policy and
continuity assertions; they do not establish numerical stability, current
compliance or the tax and expense details still absent from these sources.
